\documentclass[11pt]{article}
\usepackage{booktabs}
\usepackage{hyperref}
\usepackage{amsmath}
\title{Unknown Is Not Ineligible: Evaluating Uncertainty Preservation in\\Language-Model Benefit Decisions}
\author{Anonymous submission}
\date{}

\begin{document}
\maketitle

\begin{abstract}
Language-model agents increasingly mediate information-intensive decisions, yet a
missing fact is neither evidence for nor against eligibility. We introduce a paired,
real-record evaluation of whether agents preserve this distinction in public-benefit
reasoning. Starting from audited public-use records from the US Supplemental Nutrition
Assistance Program Quality Control database, we vary only how a decision-critical field
is presented: complete, omitted, blank, explicitly unknown, or unverified. Our primary
outcome is the Unknown-to-Denial Rate, complemented by clarification accuracy,
unsupported evidence use, and over-deferral. We compare forced binary decisions,
three-way decisions, uncertainty instructions, and an evidence-gated protocol. Results
will be inserted only after the preregistered model evaluation is completed.
\end{abstract}

\section{Introduction}
A system that lacks a required income value does not thereby know that an applicant's
income exceeds a limit. Nevertheless, binary decision interfaces and language-model
priors may transform absent evidence into an adverse conclusion. This work isolates that
failure mode using paired presentations of real audited benefit records.

Our contributions are: (1) a provenance-preserving benchmark derived from real SNAP QC
records; (2) a deterministic three-way evaluation separating ineligibility from
insufficient evidence; (3) controlled comparisons among four representations of
unresolved information; and (4) an evidence-gated protocol evaluated without training.

\section{Task and Data}
We use the FY2024 SNAP QC public-use database. It contains 44,891 active reviewed cases.
Because it is not a sample of denied applications, our estimand is model behavior under
controlled evidence removal, not the population rate of administrative denial.

\section{Experimental Design}
For each selected record, one applicable decision-critical value is shown completely,
omitted, blank, marked unknown, or marked unverified. All remaining displayed values are
identical within the matched set. Models return one of \textsc{eligible},
\textsc{ineligible}, or \textsc{insufficient information} plus requested fields and
evidence references.

\section{Metrics}
Our primary metric is
\[
\mathrm{UDR}=\frac{\#\{\text{insufficient-evidence cases predicted ineligible}\}}
{\#\{\text{insufficient-evidence cases}\}}.
\]

\section{Results}
\textbf{Intentionally blank.} No numerical result is reported before execution of the
preregistered experiment.

\section{Limitations and Ethics}
The simplified oracle does not reproduce every state rule or waiver and must not be used
for real decisions. Controlled missingness is an experimental intervention, not a claim
about natural administrative missingness. Sensitive person-level attributes are not
presented to models.

\end{document}

